III) \emph{Caso general.} --- El caso $\dim(A)=0$ es trivial, pues entonces
$A=k$; es posible, por tanto, limitarse al caso en que $\dim(A)\geqslant1$.
Se sabe (6.12.7) que $\operatorname{Spec}(A)$ contiene un abierto no vacío
$V$ cuyos puntos son todos regulares; como $\dim(A)\geqslant1$, la
intersección $V'$ de $V$ con el complementario del punto cerrado de
$\operatorname{Spec}(A)$ es un abierto no vacío y contiene, por tanto, el
punto genérico de $\operatorname{Spec}(A)$. Si demostramos que
$f^{-1}(V')\subset U_{\mathbf R}(\operatorname{Spec}(B))$,
la proposición quedará demostrada \emph{a fortiori}. En otras palabras,
basta ver que el conjunto $Z$, intersección de $f^{-1}(V')$ con el
complementario de $U_{\mathbf R}(\operatorname{Spec}(B))$, es \emph{vacío}.
Razonemos por contradicción: como, en virtud de
$(\mathrm R_{\mathrm{III}})$,
$U_{\mathbf R}(\operatorname{Spec}(B))$ es abierto en
$\operatorname{Spec}(B)$, $Z$ es localmente cerrado en
$\operatorname{Spec}(B)$; si no fuera vacío, contendría un punto $x$ tal
que $\dim(\overline{\{x\}})\leqslant1$ (5.1.10); como $f(x)\in V'$
es distinto del punto cerrado de $\operatorname{Spec}(A)$, $x$ es distinto
del punto cerrado de $\operatorname{Spec}(B)$, es decir,
$\dim(\overline{\{x\}})=1$. Vamos a demostrar que esto es imposible,
es decir, que, si $\dim(\overline{\{x\}})=1$ y $f(x)\in V'$, se cumple
$\boldsymbol R(\mathscr O_x)$. Dicho de otro modo, se trata de ver que,
si $\mathfrak q$ es un ideal primo de $B$ tal que
$\dim(B/\mathfrak q)=1$ y si $\mathfrak p=\varphi^{-1}(\mathfrak q)$
es distinto de $\mathfrak m$ y tal que $A_{\mathfrak p}$ es \emph{regular},
entonces se cumple $\boldsymbol R(B_{\mathfrak q})$. Como
$\mathfrak p\ne\mathfrak m$, $\mathfrak m(B/\mathfrak q)$ no se reduce
a 0 y, por tanto, es un ideal de definición del anillo local íntegro
$B/\mathfrak q$ de dimensión 1; por otra parte, en virtud de (i), el
cuerpo residual de $B/\mathfrak q$ es una extensión finita del cuerpo
residual $k$ de $A/\mathfrak p$, por lo que $B/\mathfrak q$ es una
$(A/\mathfrak p)$-álgebra \emph{cuasifinita} ($\mathbf{0}_{\mathrm I}$, 7.4.4). Pero
$A/\mathfrak p$ es completo y $B/\mathfrak q$ es separado para la
topología $\mathfrak m$-preádica (que coincide con su topología de anillo
local); por tanto ($\mathbf{0}_{\mathrm I}$, 7.4.1), $B/\mathfrak q$ es una
$(A/\mathfrak p)$-álgebra \emph{finita}. Además, por definición, el homomorfismo
$A/\mathfrak p\to B/\mathfrak q$ es inyectivo, por lo que (\textbf{0}, 16.1.5)
se tiene
$\dim(A/\mathfrak p)=\dim(B/\mathfrak q)=1$.
Se puede entonces aplicar a los anillos $A/\mathfrak p$ y
$B/\mathfrak pB$ el resultado de II), pues los cuerpos residuales de
estos anillos locales son, respectivamente, los de $A$ y $B$, y
$B/\mathfrak pB$ es un $(A/\mathfrak p)$-módulo plano; además, se tiene
$(B/\mathfrak pB)\otimes_{A/\mathfrak p}k=B\otimes_A k=B/\mathfrak mB$.
Por consiguiente, los anillos locales de
$\operatorname{Spec}((B/\mathfrak pB)
\otimes_{A/\mathfrak p}\boldsymbol{k}(\mathfrak p))$
satisfacen $\boldsymbol R$. Ahora bien,
$B_{\mathfrak q}/\mathfrak pB_{\mathfrak q}$ es uno de los anillos
locales de $\operatorname{Spec}((B/\mathfrak pB)
\otimes_{A/\mathfrak p}\boldsymbol{k}(\mathfrak p))$. Además, $B_{\mathfrak q}$ es
un $A_{\mathfrak p}$-módulo plano y $A_{\mathfrak p}$ es \emph{regular}.
Se tiene, pues, el siguiente lema:
